% ---------------------------------------------------------------------------
%  00-how-to-read.tex
% ---------------------------------------------------------------------------
\section*{How to read this}
\addcontentsline{toc}{section}{How to read this}

Two audiences. Someone deciding whether the idea is worth a paper should read
\cref{sec:problem} (what the measurement says the problem actually is),
\cref{sec:theory} (the theorem), and \cref{sec:eval} (what it buys). Someone
maintaining the code should read \cref{sec:algorithm} and keep
\cref{sec:negative} nearby, because most of what looks like an optimisation here
is not one, and that section says which.

Every mechanism is traceable to one place in the source.

\begin{center}\small
\begin{tabular}{@{}lll@{}}
\toprule
Mechanism & Where & Section \\
\midrule
partition a node into components   & \code{components()}            & \cref{sec:algorithm} \\
decide a node by its components    & \code{decide\_by\_components()} & \cref{sec:algorithm} \\
hook into the search               & \code{try\_split} in \code{decide\_scc} & \cref{sec:algorithm} \\
atom sets as bitmasks              & \code{atoms\_of()}             & \cref{sec:algorithm} \\
clock ownership                    & \code{clocks\_of()}            & \cref{sec:algorithm} \\
shapes that never split            & \code{nosplit\_}               & \cref{sec:algorithm} \\
root-level split (static form)     & \code{Tableau::check}          & \cref{sec:eval} \\
\bottomrule
\end{tabular}
\end{center}

\noindent
Both forms are behind flags and default off: \code{TCTL\_COMPONENTS=1} for the
per-node form, \code{TCTL\_DECOMPOSE=1} for the root-level one. With neither set
the search is exactly what it was.

\begin{intuition}
The tableau's cost has two completely different causes, and they need different
cures. One is that a node carries many clock regions; the other is that it
carries many independent obligations. Everything anyone would normally try
attacks the first. The measurement in \cref{sec:problem} says the expensive
cases are the second, and the second has a clean fix.
\end{intuition}
